% !TeX root = ../泛函分析笔记.tex
% 由主文档通过 \input 载入；请编译主文档。
\setcounter{chapter}{2} % 保留原书章号。
\chapter{赋范线性空间与 Banach 空间}
\label{chap:fa-banach}

在向量空间中，极限还必须与加法、数乘相容。范数同时描述向量的大小和向量间的距离，是建立无限维线性分析的基本结构。

\section{范数、半范数与等价范数}
\begin{definition}
$\mathbb K$ 上线性空间 $X$ 的\emph{范数（norm）}是映射 $\|\cdot\|:X\to[0,\infty)$，满足
\[
 \|x\|=0\iff x=0,\qquad \|\alpha x\|=|\alpha|\|x\|,\qquad
 \|x+y\|\leq\|x\|+\|y\|.
\]
不要求第一项的正定性时称为半范数（seminorm）。范数诱导度量 $d(x,y)=\|x-y\|$；此度量完备时，$X$ 称为\emph{Banach 空间（Banach space）}。
\end{definition}
反三角不等式 $|\|x\|-\|y\||\leq\|x-y\|$ 表明范数连续。若 $x_n\to x$、$y_n\to y$、$\alpha_n\to\alpha$，则三角不等式以及 $(x_n)$ 的有界性给出 $x_n+y_n\to x+y$ 和 $\alpha_nx_n\to\alpha x$。

同一空间的两个范数 $\|\cdot\|_a,\|\cdot\|_b$ 称为等价，若存在 $c,C>0$ 使
\[
 c\|x\|_a\leq\|x\|_b\leq C\|x\|_a\quad(x\in X).
\]
两侧估计表明等价范数有相同的收敛序列、Cauchy 序列和开集，因而完备性也相同。反过来，若两个范数给出相同的拓扑，则两个恒等映射在零点连续，利用线性缩放就得到上述估计。

\section{数列空间与连续函数空间}
有限维中常见的三个范数是 $\|x\|_1=\sum_{j=1}^d|x_j|$、$\|x\|_2=(\sum_{j=1}^d|x_j|^2)^{1/2}$ 和 $\|x\|_\infty=\max_j|x_j|$。例如 $x=(3,4)$ 的三个范数分别是 $7,5,4$。它们衡量大小的方式不同，但满足 $\|x\|_\infty\leq\|x\|_2\leq\|x\|_1\leq d\|x\|_\infty$，因此给出相同的收敛。把有限求和改成无限求和时，首先要筛选使求和有意义的数列，这就产生下面的空间。

对 $1\leq p<\infty$ 定义
\[
 \ell^p=\left\{x=(x_j)_{j\geq1}:\sum_{j=1}^\infty|x_j|^p<\infty\right\},
 \qquad \|x\|_p=\left(\sum_{j=1}^\infty|x_j|^p\right)^{1/p}.
\]
另令 $\ell^\infty=\{x:\sup_j|x_j|<\infty\}$，$\|x\|_\infty=\sup_j|x_j|$；$c_0$ 为其中趋于零的数列，$c_{00}$ 为仅有限个坐标非零的数列。$\ell^p$ 的三角不等式是 Minkowski 不等式，见附录\ref{app:fa-measure}；计数测度将积分形式化为求和形式。

\begin{example}[一个空间中的点本身也可以是数列]
数列 $x=(1,1/2,1/3,\ldots)$ 属于 $\ell^2$ 与 $c_0$，但不属于 $\ell^1$；常数列 $\boldsymbol 1=(1,1,\ldots)$ 属于 $\ell^\infty$，却不属于 $c_0$。一般地，对 $a>0$，$(j^{-a})\in\ell^p$ 当且仅当 $ap>1$。
\end{example}
\begin{solve}
这些判断分别归结为级数 $\sum j^{-2}$ 收敛、调和级数发散，以及一般正项级数 $\sum j^{-ap}$ 的判别。注意区分两个指标：$x_j$ 是一个向量的第 $j$ 个坐标，而 $x^{(n)}$ 是空间中第 $n$ 个向量。$x^{(n)}\to x$ 是整条数列作为向量趋近另一条数列，不等于仅检查每个固定坐标。
\end{solve}

\begin{proposition}\label{prop:fa-sequence-complete}
$\ell^p$（$1\leq p\leq\infty$）和配备上确界范数的 $c_0$ 都是 Banach 空间。若 $K$ 为非空紧度量空间，则 $C(K)$ 在一致范数 $\|f\|_\infty=\max_{t\in K}|f(t)|$ 下也是 Banach 空间。
\end{proposition}
\begin{proof}
先设 $p<\infty$，$(x^{(n)})$ 为 $\ell^p$ 中的 Cauchy 序列。由 $|x_j^{(n)}-x_j^{(m)}|\leq\|x^{(n)}-x^{(m)}\|_p$，每个坐标收敛，记极限为 $x_j$。若 $n,m\geq N$ 时范数差小于 $\varepsilon$，则先对有限个坐标求和并令 $m\to\infty$，再令坐标数趋于无穷，得到
\[
 \sum_j|x_j^{(n)}-x_j|^p\leq\varepsilon^p\quad(n\geq N).
\]
因此 $x-x^{(N)}\in\ell^p$，从而 $x\in\ell^p$ 且 $x^{(n)}\to x$。$p=\infty$ 时，逐坐标取极限后保留同一上界即可。$c_0$ 在 $\ell^\infty$ 中闭：先用一个趋零数列一致逼近极限数列，再控制其尾部。$C(K)$ 的证明与上一章 $C([0,1])$ 的一致极限证明相同。
\end{proof}
截断前 $N$ 个坐标表明，$c_{00}$ 在每个 $\ell^p$（$p<\infty$）和 $c_0$ 中稠密。用实部、虚部为有理数的有限支撑数列逼近，可见这些空间可分；$\ell^\infty$ 不可分，因为所有零一数列构成不可数集，任意两项的距离为一。

这里截断误差可以直接计算。令 $P_Nx=(x_1,\ldots,x_N,0,\ldots)$，则 $p<\infty$ 时 $\|x-P_Nx\|_p^p=\sum_{j>N}|x_j|^p\to0$；在 $c_0$ 中误差为 $\sup_{j>N}|x_j|\to0$。但对 $\boldsymbol1\in\ell^\infty$，截断误差始终等于一。为了补全不可分性的论证，若存在可数稠密集，每个零一数列附近半径 $1/3$ 的开球都应含有稠密集的一点；这些球两两不交，这将需要不可数多个不同的稠密集元素，矛盾。

\section{可测函数空间与积分范数}
给定测度空间 $(\Omega,\mathcal A,\mu)$，$L^p(\mu)$ 的元素是几乎处处相等的可测函数的等价类。对 $1\leq p<\infty$ 令
\[
 \|f\|_p=\left(\int_\Omega|f|^p\,\mathrm d\mu\right)^{1/p},
 \qquad
 \|f\|_\infty=\inf\{C\geq0:|f|\leq C\ \text{几乎处处}\}.
\]
只保留范数有限的等价类。取等价类是必要的：在原函数集合上，积分范数为零并不意味着处处为零。

比如在 Lebesgue 测度下，恒零函数与只在 $t=1/2$ 处取一的函数代表同一个 $L^p(0,1)$ 元素。因此对一个抽象的 $L^p$ 元素说“它在 $1/2$ 处的值”，通常没有定义；但它的积分和范数都有定义。$\|f\|_\infty$ 中的本质上确界（essential supremum）也忽略零测集上的异常值，与对每一点取上确界不同。

对 $f(t)=t^{-a}$、$a>0$，有 $f\in L^p(0,1)$ 当且仅当 $ap<1$，因为 $\int_0^1t^{-ap}\,\mathrm dt$ 的收敛条件如此。这个条件与数列 $(j^{-a})$ 的条件方向相反：积分例子检查零点附近的奇性，数列例子检查无穷远的尾部。在有限测度空间上，若 $1\leq p<q\leq\infty$，Hölder 不等式给出
\[
 \|f\|_p\leq\mu(\Omega)^{1/p-1/q}\|f\|_q
 \quad(0<\mu(\Omega)<\infty),
\]
故 $L^q\subset L^p$；数列空间则是 $\ell^p\subset\ell^q$。使用包含关系前，必须先说明测度空间，不能只记一个箭头方向。

\begin{theorem}\label{thm:fa-lp-complete}
对任意测度空间及 $1\leq p\leq\infty$，$L^p(\mu)$ 是 Banach 空间。
\end{theorem}
\begin{proof}
设 $p<\infty$。从 Cauchy 序列 $(f_n)$ 选取子列，使 $\|f_{n_{k+1}}-f_{n_k}\|_p\leq2^{-k}$。令 $g_k=|f_{n_{k+1}}-f_{n_k}|$。Minkowski 不等式和单调收敛定理给出 $\|\sum_kg_k\|_p\leq\sum_k2^{-k}<\infty$，故该和几乎处处有限。于是函数级数绝对收敛，定义
\[
 f=f_{n_1}+\sum_{k=1}^{\infty}(f_{n_{k+1}}-f_{n_k})
 \quad\text{几乎处处}.
\]
尾项的 $L^p$ 范数不超过相应几何级数尾和，故子列在范数下趋于 $f\in L^p$，原 Cauchy 序列也趋于 $f$。当 $p=\infty$ 时，去掉控制各对差的可数个零测集后，序列在剩余集合上一致 Cauchy，逐点极限便给出 $L^\infty$ 极限。
\end{proof}
范数收敛不要求原序列逐点收敛。例如在 $[0,1)$ 上依次列出各级二进区间的示性函数，其 $L^p$ 范数在 $p<\infty$ 时趋于零，但每个点都被无穷多个区间包含，也被无穷多个区间排除。

\begin{example}[为什么要从连续函数走向可积函数]\label{ex:fa-incomplete-integral}
对 $n\geq3$，令 $f_n$ 在 $[0,1/2]$ 上为零，在 $[1/2,1/2+1/n]$ 上线性地从零升到一，在剩余区间上为一。则 $(f_n)$ 在积分范数下为 Cauchy 序列，但不在 $C([0,1])$ 中收敛。
\end{example}
\begin{solve}
令 $g=\mathbf1_{(1/2,1]}$。两者只在宽度 $1/n$ 的过渡区间有差别，三角形面积给出 $\|f_n-g\|_1=1/(2n)$，所以 $\|f_n-f_m\|_1\leq1/(2n)+1/(2m)$。若存在连续 $f$ 使 $\|f_n-f\|_1\to0$，则 $\|f-g\|_1=0$。连续性迫使 $f$ 在 $(0,1/2)$ 上恒为零、在 $(1/2,1)$ 上恒为一：否则某点附近有正长度区间上的差不为零，积分不可能为零。这与 $f$ 在 $1/2$ 连续矛盾。新空间 $L^1$ 正好容纳了这个极限，而一致范数下的连续函数空间本来就是完备的。
\end{solve}

\section{子空间、商空间与级数判别}
Banach 空间的线性子空间完备当且仅当它闭，这是命题\ref{prop:fa-closed-complete}的直接应用。若 $M\subset X$ 为闭线性子空间，在商空间 $X/M$ 上定义
\begin{equation}
 \|x+M\|_{X/M}=\inf_{m\in M}\|x-m\|_X.
 \label{eq:fa-quotient-norm}
\end{equation}
右侧也记为 $\operatorname{dist}(x,M)$，即点到集合的距离。更换代表元不会改变右侧；齐次性和三角不等式由取近似极小元得到。范数为零等价于 $x\in\overline M=M$，所以闭性保证正定性。

商空间把相差一个 $M$ 中元素的向量视为相同。其零元素是整个集合 $M$，而不是只含原点的集合。例如 $X=\mathbb R^2$、$M=\mathbb R\times\{0\}$ 时，$(a,b)+M$ 是高度为 $b$ 的水平直线；在欧氏范数下，$\|(a,b)+M\|=|b|$，故商空间等距对应于纵坐标轴。函数空间中，对常数函数构成的子空间取商，就把“相差一个常数”的函数视为同一对象。微分算子无法区分这些函数，因而商空间自然适合消除这类不唯一性。

商范数中的下确界未必有最小元。证明三角不等式时只需对任意 $\varepsilon>0$，分别选 $m_1,m_2\in M$，使 $\|x-m_1\|<\|x+M\|+\varepsilon$、$\|y-m_2\|<\|y+M\|+\varepsilon$。用 $m_1+m_2$ 作为 $x+y$ 的一个候选代表，再令 $\varepsilon\downarrow0$ 即可；不应事先假定存在最近点。

\begin{proposition}\label{prop:fa-series-complete}
赋范空间完备当且仅当每个满足 $\sum_n\|x_n\|<\infty$ 的级数 $\sum_nx_n$ 都收敛。Banach 空间对闭线性子空间取商后仍为 Banach 空间。
\end{proposition}
\begin{proof}
完备空间中，绝对可和级数的部分和为 Cauchy 序列。反之，对任意 Cauchy 序列 $(y_n)$ 选子列使 $\|y_{n_{k+1}}-y_{n_k}\|\leq2^{-k}$；差分级数收敛，所以子列及原序列均收敛。

设 $\sum_n\|\xi_n\|_{X/M}<\infty$。对每个商类 $\xi_n$ 选择代表元 $x_n$ 使 $\|x_n\|\leq\|\xi_n\|_{X/M}+2^{-n}$。则 $\sum_nx_n$ 在 $X$ 中收敛，连续商映射把其和映为 $\sum_n\xi_n$ 的极限。由刚证的级数判别，$X/M$ 完备。
\end{proof}
赋范空间的度量完备化可通过代表序列逐项相加、数乘，赋予唯一延续原运算的 Banach 空间结构；范数定义为原序列范数的极限。

级数判别中的“绝对”不能省略，也不表示所有收敛级数都绝对收敛。在 $\ell^2$ 中，$\sum_{n\geq1}e_n/n$ 收敛到 $(1/n)$，因为尾项范数平方为 $\sum_{n>N}1/n^2\to0$；但各项范数之和为 $\sum1/n=\infty$。完备性保证绝对可和是收敛的充分条件，并未把它变成必要条件。

\section{有限维定理与 Riesz 引理}
\begin{theorem}\label{thm:fa-finite-dimension}
有限维线性空间上的所有范数等价；有限维赋范空间完备，其闭有界集紧；任意赋范空间的有限维线性子空间闭。
\end{theorem}
\begin{proof}
取基 $e_1,\ldots,e_d$，写 $x=\sum_ja_je_j$。三角不等式给出 $\|x\|\leq C\|a\|_2$，故 $a\mapsto\|\sum_ja_je_j\|$ 在欧氏单位球面上连续且正，存在正最小值 $c$。由齐次性得 $c\|a\|_2\leq\|x\|$，从而任意范数等价于坐标欧氏范数。完备与紧性由欧氏空间转移；有限维子空间完备，因而在环境空间中闭。
\end{proof}
\begin{lemma}[Riesz 引理]\label{lem:fa-riesz}
若 $M$ 是赋范空间 $X$ 的真闭线性子空间，$0<\alpha<1$，则存在单位向量 $x$ 满足 $\inf_{m\in M}\|x-m\|>\alpha$。
\end{lemma}
\begin{proof}
取 $y\notin M$，令 $d=\inf_{m\in M}\|y-m\|>0$，选 $m_0\in M$ 使 $\|y-m_0\|<d/\alpha$。令 $x=(y-m_0)/\|y-m_0\|$，则对任意 $m\in M$，$\|x-m\|\geq d/\|y-m_0\|>\alpha$。
\end{proof}
\begin{corollary}\label{cor:fa-unit-ball}
赋范空间的闭单位球紧当且仅当空间有限维。
\end{corollary}
\begin{proof}
无限维时，对逐步生成的有限维子空间反复使用引理\ref{lem:fa-riesz}，得到单位向量 $x_n$，任意两项距离大于 $1/2$，因而闭单位球不列紧。有限维方向由定理\ref{thm:fa-finite-dimension}得到。
\end{proof}

Riesz 引理在没有内积的空间中起到寻找“新方向”的作用。在 Hilbert 空间里可以选择与已有有限维空间正交的单位向量，其距离恰为一；一般赋范空间没有角度，却仍能找到与已有空间保持固定距离的单位向量。由此可见，无限维单位球不紧并非某个特殊范数造成的偶然现象，而是任意范数都会遇到的结构性障碍。

\section{例题与习题}
\begin{example}[不同范数的作用]
$C([0,1])$ 上的 $\|f\|_1=\int_0^1|f|$ 与 $\|f\|_\infty$ 不等价。
\end{example}
\begin{solve}
虽有 $\|f\|_1\leq\|f\|_\infty$，但取 $f_n(t)=\max(1-nt,0)$，则 $\|f_n\|_\infty=1$、$\|f_n\|_1=1/(2n)$，反向统一估计不可能成立。这说明在无限维空间中，改变范数可能改变收敛与连续性。
\end{solve}
\begin{enumerate}
 \item 证明 $c_{00}$ 在 $\ell^\infty$ 中的闭包为 $c_0$。
 \item\label{exer:fa-lp-inclusion} 证明 $\ell^p\subset\ell^q$ 且 $\|x\|_q\leq\|x\|_p$，其中 $1\leq p<q\leq\infty$。
 \item 在 $\mathbb R^2$ 的欧氏范数下，计算对 $x$ 轴取商后的范数。
 \item 对例\ref{ex:fa-incomplete-integral}中的函数列，计算 $\|f_n-g\|_p$（$1\leq p<\infty$），说明为何同一反例适用于每个有限积分范数。
\end{enumerate}

